The Adams (smallest divisors) apportionment method assigns each state
$s_i = \lceil p_i/d \rceil$ seats for a common divisor $d$ chosen to
produce a total of $H$ seats.
Ceiling rounding --- always rounding up --- means every state receives
at least as many seats as its exact proportional quota demands,
making Adams the most small-state-favoring of the five major apportionment
methods.

We prove that Adams maximises the minimum per-capita representation across
all states: $\max_d \min_i (p_i / \lceil p_i/d \rceil)$.
Equivalently, Adams guarantees that no state's per-capita representation
falls below the level achieved by its upper quota entitlement.
This property makes Adams the natural benchmark for maximum small-state
protection in apportionment reform debates.

We prove paradox immunity for the Adams method using the same
divisor-method monotonicity argument that applies to Huntington-Hill
and Webster: because the seat allocation is determined by a priority
queue that only adds seats, increasing the House size or increasing a
state's population can never reduce any state's seat count.

Using 2020 Census data, Adams gives Rhode Island, Montana, and Alaska
each one additional seat relative to Huntington-Hill,
while large states (California, Texas, Florida) correspondingly receive
fewer seats to maintain the 435-seat total.
We document that Adams has not been used for federal apportionment
in the United States but is used in some European subnational
proportional representation systems.
The \textsc{bisect-apportion} crate implements Adams with the same
priority-queue architecture as the other divisor methods.
